
avec \(G\) un \((b_1,\ldots,b_r)\)-faisceau et \(\operatorname{dim{.}supp}(G)=s-1\). Alors
\[
h^0(F(n))-h^0(F(n-1))\le h^0(G(n)),
\]
mais
\[
h^0(G(n))\le\sum_{i=0}^{s-1}b_{i+1}\binom{n+i}{i}
\]
par récurrence sur \(s\) et \(h^0(F(-1))\le b_0\); d'où (i) par récurrence sur \(n\ge -1\). Enfin, par 1.7 et récurrence sur \(s\), \(a_s\le b_s\) et \(G\) est un \((b_i,\ldots,b_{s-1},a_s)\)-faisceau; d'où (ii).

\textbf{Définition 1.9.} \emph{On appelle \emph{\((b)\)-polynômes} les polynômes à coefficients rationnels \(\ge0\) suivants, définis par récurrence}~:
\[
\begin{cases}
P_{-1}=0\\
P_r(X_0,\ldots,X_r)=P_{r-1}(X_1,\ldots,X_r)
+\displaystyle\sum_{i=0}^{r}X_i
\binom{P_{r-1}(X_1,\ldots,X_r)-1+i}{i}.
\end{cases}
\]

\textbf{Remarque 1.10.} Pour les \((b)\)-polynômes, on voit tout de suite par récurrence sur \(r\ge t\) que
\[
P_r(X_0,\ldots,X_t,0,\ldots,0)=P_t(X_0,\ldots,X_t).
\]

\textbf{Théorème 1.11.} \emph{(Le théorème principal sur les \((b)\)-faisceaux). Soient \(F\) un \((b)\)-faisceau sur \(X\) et \(\chi(F(n))=\sum_{i=0}^{r}a_i\binom{n+i}{i}\) le polynôme de Hilbert de \(F\); soient \((c)=(c_0,\ldots,c_r)\) une suite d'entiers \(\ge (b)-(a)\) et \(m=P_r(c)\) l'entier donné par le \(r\)-ième \((b)\)-polynôme (1.9). Alors \(m\ge0\) et \(F\) est \(m\)-régulier. En particulier, si \(s=\operatorname{dim{.}supp}(F)\), alors \(F\) est \(P_{s-1}(c_0,\ldots,c_{s-1})\)-régulier.}

